\section{Characteristic-zero parameters and local control}
\label{sec:parameters}

We attach Galois representations to the genuine cusp systems of
\cref{prop:cusp-congruence}, retain the local information supplied
by the prescribed operators, and then remove the nilpotents in the
operator order with a uniform loss of congruence precision.
The local information takes the form of invariant subspaces and
polynomial identities in inertia matrices.  These identities will
also impose the local conditions on the extensions extracted in
\cref{sec:extensions}.

\subsection{Simultaneous automorphic transfer}
\label{subsec:simultaneous-transfer}

Put $G=\U(B\oplus\mathbb H)$, with the Hermitian space and conventions of
\cref{sec:theta}.  Thus $G(F_v)\simeq\U(3,1)$ at each real place,
$G$ is quasi-split at every finite place, and its proper $F$-Levi subgroup
has real points $\C^\times\times\U(2)$ at either real place.
Let $\pi$ be a cuspidal automorphic representation occurring in
\cref{prop:cusp-congruence}.  Its infinite components are holomorphic
discrete series of lowest compact type $\lambda'$ from
\cref{eq:shifted-weight}.  At the finite places nonsplit in $L$, its
components are hyperspecial spherical.

For a unitary cuspidal representation $\tau$ of $\GL_d(\A_L)$, write
$\tau[a]$ for its Arthur block of length $a$ and put
\[
 S_a=\{(a-1)/2,(a-3)/2,\ldots,-(a-1)/2\}.
\]
The local standard representation attached to this block is the
Langlands quotient of its shifted cuspidal data.  In particular, for a
finite place $w$ of $L$ its Weil--Deligne parameter is
\begin{equation}
 \operatorname{rec}_w(\tau_w[a])
   =\bigoplus_{s\in S_a}
        \operatorname{rec}_w(\tau_w)\otimes|\cdot|_w^s.
 \label{eq:arthur-block-wd}
\end{equation}
The monodromy on the right is the direct sum of the monodromies of the
shifted copies.  This is the Arthur-block convention, as opposed to a
new Steinberg block.  Local Langlands and the classification of the
unitary dual of general linear groups give
\cref{eq:arthur-block-wd}; normalized products of the resulting unitary
representations are irreducible.  At split places these character
identities are the linear-group case of endoscopic transfer, also
treated in \citet[Section~1.6.3]{KMSW2014}.

We obtain one formal sum of these blocks, with its full infinite
character and finite local parameters, from
\citet[Corollary~5.3]{Labesse2011CM}.  The needed sign property of the
individual holomorphic discrete series is verified below at each real
place.  Labesse's simple trace-formula argument applies because
$[F:\Q]=2$ and does not use the weighted fundamental lemma
\citep[Sections~4.8 and~5.1]{Labesse2011CM}.
In \cref{eq:shifted-weight} we take $H_i\geq3$, as we may while
preserving the specified divisibility and congruences.

\begin{lemma}[Weak transfer with full infinite characters]
\label{lem:weak-transfer}
There is a conjugate-dual formal parameter
\[
 \Psi=\boxplus_\beta\tau_\beta[a_\beta],
 \qquad \sum_\beta d_\beta a_\beta=4,
\]
with unitary cuspidal factors, the standard transferred full
infinitesimal character of $\pi$ at both real places, and the
transferred Satake parameter at every finite place outside a finite
set $S$ whose finite places split in $L/F$.
Every cuspidal factor is unramified at every nonsplit finite place.
\end{lemma}
\begin{proof}
First consider one real factor $\U(3,1)$, and omit the subscript $i$.
The Harish--Chandra coordinates of the selected holomorphic discrete
series $D$ are
\[
 k=(m+H+\tfrac12,\ m+H-\tfrac12,\ H-\tfrac12,\ \tfrac12-H).
\]
Their successive gaps are $1,m,2H-1$.  They are regular integral,
and the unique negative-sign coordinate of $D$ is $k_4$.
Choose its individual pseudocoefficient $f_D$, with trace one on $D$
and trace zero on every other basic standard module
\citep[Proposition~4 and its corollary, pp.~212--213]{ClozelDelorme1990}.
In particular it annihilates all properly parabolically induced
standard modules.  Standard-module expansions preserve infinitesimal
character, so an irreducible representation with nonzero trace against
$f_D$ must have infinitesimal character $k$.

Every unitary representation with this regular integral character is
an admissible cohomologically induced module $A_{\mathfrak q}(\mu)$
\citep[Theorem~1.2]{SalamancaRiba1988}.  That theorem applies to
$\operatorname{SU}(3,1)$; adjoining the compact central character gives the same
assertion for $\U(3,1)$.  Write the theta-stable Levi of
$\mathfrak q$ as ordered signature blocks
$\prod_j\U(p_j,q_j)$, where $\sum p_j=3$ and $\sum q_j=1$.
The inducing character is constant on each Levi block and dominant
across the nilradical roots.  Hence the infinitesimal coordinates
in each block form a consecutive string, with successive difference
one, and the blocks occur in their decreasing order in $k$.
There is exactly one block with $q_j=1$.

If this Levi is noncompact, its block containing the negative
coordinate has size at least two.  In the standard-module resolution
of $A_{\mathfrak q}(\mu)$, the discrete-series terms arise from
positive systems extending the nilradical roots
\citep[Theorem~8.2 and Lemma~8.8]{AdamsJohnson1987}.
They may reorder coordinates within a Levi block, but not between
blocks.  Their negative coordinate must therefore belong to that
same block.  For $D$ to occur, the block would have to contain $k_4$.
A consecutive string of size at least two containing $k_4$ also
contains $k_3$ and forces $k_3-k_4=1$, contrary to
$k_3-k_4=2H-1\geq5$.  Thus the coefficient of $D$ in this resolution
is zero.  All other terms are annihilated by $f_D$.
If the Levi is compact, the induced module is itself discrete
series, and its trace is one precisely when it is $D$.
We have proved that $\operatorname{tr}X(f_D)$ is zero or one for
every irreducible unitary $X$, with value one only for $X=D$.
This argument treats each real factor separately.  Their product
pseudocoefficient therefore satisfies the sign property $(*)$ of
\citet[Section~5.1]{Labesse2011CM}, including competitors with
different types at the two real places.  No lower bound on the
first gap beyond its actual value one is used.

Apply \citet[Corollary~5.3]{Labesse2011CM} to $\pi$.
It gives one automorphic transfer induced from conjugate-dual
discrete general-linear representations.  The discrete-spectrum
classification of \citet{MoeglinWaldspurger1989} expresses its
factors as Arthur blocks and gives the formal parameter
\[
 \Psi=\boxplus_\beta\tau_\beta[a_\beta],
 \qquad \sum_\beta d_\beta a_\beta=4,
\]
with unitary cuspidal factors.  The same corollary gives local
base change at every split finite place and at every unramified
place.  It also retains the full infinitesimal character at both
real places: its proof uses the coefficient systems for the
complex group and compatibility with infinitesimal-character
multipliers, not just the Casimir eigenvalue.  We do not invoke
its additional conclusion requiring cuspidality of the transfer.

Choose $S$ to contain the infinite places and the finite ramification
of $\pi$ and the chosen standard embedding.  All its finite places
can be chosen split in $L/F$.  At a nonsplit finite place, the
transferred representation is spherical, so its Weil--Deligne
parameter has trivial inertia and zero monodromy.
Formula~\eqref{eq:arthur-block-wd} and compatibility of local
Langlands with the normalized product express this parameter as
the direct sum of the shifted parameters of $\tau_{\beta,w}$.
Each summand therefore has trivial inertia and zero monodromy,
so each $\tau_{\beta,w}$ is unramified.
The corollary's stronger matching at every split
place concerns this same $\Psi$ and will also be used below.
\end{proof}

\begin{lemma}[Algebraicity of the cuspidal factors]
\label{lem:transfer-algebraicity}
Every cuspidal factor in \cref{lem:weak-transfer} is
conjugate-self-dual and tempered at infinity.  After a unitary
parity-correcting character twist it is regular algebraic.
The formal parameter is elliptic: its conjugate-self-dual blocks
occur without repetition.
\end{lemma}

\begin{proof}
At a distinguished complex embedding the first infinitesimal
coordinates are
\begin{equation}
 (k_1,k_2,k_3,k_4)
 =\bigl(\lambda'_1+\tfrac12,\lambda'_2-\tfrac12,
        \lambda'_3-\tfrac32,\lambda'_4+\tfrac32\bigr).
 \label{eq:transfer-infinitesimal}
\end{equation}
Indeed the half-sums of positive compact and noncompact roots are
$\rho_c=(1,0,-1;0)$ and
$\rho_n=(\tfrac12,\tfrac12,\tfrac12;-\tfrac32)$, and the
holomorphic discrete-series formula is $\lambda'+\rho_c-\rho_n$.
The coordinates are distinct and belong to $\Z+\tfrac32$; the
coordinates at the conjugate embedding form the multiset
$\{-k_1,-k_2,-k_3,-k_4\}$.

Consider an archimedean character $z^u\bar z^v$ in a cuspidal factor
$\tau$ before inserting a shift $s\in S_a$.  The two coordinates
$u+s$ and $v+s$ occur in the two prescribed infinitesimal
characters.  Hence $u,v$ are real and $u+v\in\Z$.
A unitary generic representation of $\GL_d(\C)$ has
$|\operatorname{Re}(u+v)|<1$ for every inducing character.
Cuspidality supplies genericity, so $u+v=0$.  This proves
temperedness at infinity.

The archimedean parameters of $\tau$ and $\tau^{c,\vee}$ are now
identical.  A pair of distinct conjugate-dual cuspidal factors in
$\Psi$ would consequently repeat the first infinitesimal
coordinates, contrary to regularity.  Thus all factors are
conjugate-self-dual.  Repeated copies of an identical Arthur block
are excluded for the same reason.  In particular the resulting
unitary parameter is elliptic.

Choose a unitary Hecke character $\chi$ of $L$ with
$\chi^c=\chi^{-1}$, with infinity type $z/|z|$ at both
distinguished embeddings, and with conductor supported on split
auxiliary primes.  Its existence is the character construction of
\cref{sec:conventions}, with the following parity prescription.
Require $\chi|_{\A_F^\times}=\omega_{L/F}$, the quadratic
character attached to $L/F$.  At infinity this agrees with the
restriction of $z/|z|$ to $\R^\times$.  At nonsplit finite places
use unramified extensions of $\omega_{L/F,v}$, and use a sufficiently
deep conductor at split auxiliary primes to remove the remaining
ray-class unit obstruction.  The resulting Hecke character satisfies
\[
 \chi(x)\chi(x^c)=\chi(\Norm_{L/F}x)
      =\omega_{L/F}(\Norm_{L/F}x)=1,
\]
which gives the required global identity $\chi^c=\chi^{-1}$.
For a factor $\tau$ of rank $d$ and block
length $a$, choose $\delta\in\{0,1\}$ with
\[
 \delta\equiv a+d-1\pmod2.
\]
If $u$ is one of its first infinitesimal coordinates, then
$u\in\Z+(4-a)/2$, and therefore
$u+\delta/2\in\Z+(d-1)/2$.
Thus $\Pi=\tau\otimes\chi^\delta$ is regular algebraic in the
cohomological normalization for $\GL_d$.
Regularity follows because its shifted coordinates form a subset
of the distinct coordinates in \cref{eq:transfer-infinitesimal}.
The equality $\Pi^c\simeq\Pi^\vee$ follows from those for
$\tau$ and $\chi$.
\end{proof}

\begin{proposition}[Simultaneous transfer]
\label{prop:simultaneous-transfer}
There is a formal parameter
\[
 \Psi=\boxplus_\beta\tau_\beta[a_\beta],
 \qquad \sum_\beta d_\beta a_\beta=4,
\]
with unitary cuspidal factors on $\GL_{d_\beta}(\A_L)$, such that:
\begin{enumerate}[label=\textup{(\roman*)}]
 \item its archimedean infinitesimal character is the standard transfer
 of that of $\pi$;
 \item at every split finite place $v=ww^c$, the representation
 $\pi_v$, viewed on $\GL_4(F_v)$ using $w$, is the representation
 attached to $\Psi_w$;
 \item at every nonsplit finite place, the standard local parameter
 is unramified and has zero monodromy.
\end{enumerate}
\end{proposition}

\begin{proof}
Take the single weak parameter supplied by \cref{lem:weak-transfer}.
Its archimedean character is the required one, and its unramified
cuspidal factors at nonsplit places give assertion \textup{(iii)}
by \cref{eq:arthur-block-wd}.
The application of Labesse's corollary in that lemma already gives
local base change at every split finite place.  Under the chosen
identification $G(F_v)=\GL_4(F_v)$ this is exactly assertion
\textup{(ii)}, including the places above $p$, the bad places, and
the moving levels.  All three assertions thus concern the same
parameter.  This transfer argument places no restriction on the
coefficient prime $p$.
\end{proof}

\subsection{Galois representations and their weights}

We now attach Galois representations to the algebraic twists of the
cuspidal factors in the simultaneous parameter.  The resulting
four-dimensional representation retains the direct-sum decomposition
over the shifted copies of every Arthur block.  This decomposition
will matter when we control its actual monodromy.

We use the following established Galois input, in geometric
reciprocity and with de Rham weight $1$ for $\eps^{-1}$.
\begin{theorem}[Regular algebraic conjugate-self-dual Galois input]
\label{thm:racsd-galois-input}
Let $\Pi$ be a regular algebraic conjugate-self-dual cuspidal
representation of $\GL_d(\A_L)$ and fix
$\iota:\overline{\Q}_p\simeq\C$.  There is a continuous
semisimple representation $r_p(\Pi)$ of $G_L$, with
\[
 r_p(\Pi)^c\simeq r_p(\Pi)^\vee\eps^{1-d}.
\]
It is de Rham above $p$.  If its first infinitesimal coordinates
at an embedding are $b_1,\ldots,b_d$, its de Rham weights there
are $(d-1)/2-b_j$.  At every finite place $w$,
\begin{equation}
 \iota\WD\bigl(r_p(\Pi)|_{G_{L_w}}\bigr)^{\mathrm{F\text{-}ss}}
 \preceq
 \operatorname{rec}_w
     \bigl(\Pi_w\otimes|\det|_w^{(1-d)/2}\bigr).
 \label{eq:racsd-compatibility}
\end{equation}
Here $\preceq$ means equality after semisimplifying the Weil
representation and domination of monodromy partitions on each
inertial class.  At spherical places away from $p$ the
representation is unramified, and at spherical places above $p$
it is crystalline.
\end{theorem}

Existence, polarization and the de Rham assertions are
\citet[Theorem 1.2]{BGHT2011}; compatibility away from $p$
in this strength is
\citet[Theorem 3.2.3]{ChenevierHarris2013}.
The assertion above $p$, including the monodromy bound for
arbitrary regular algebraic weight, is
\citet[Theorem 1.1]{BLGGT2014}.
The rank-one case is algebraic global class field theory.

\begin{proposition}[Galois representations of the cusp systems]
\label{prop:galois-package}
For every characteristic-zero cusp eigensystem in
\cref{prop:cusp-congruence}, there is a continuous semisimple
representation
\[
 \rho_\pi:G_L\longrightarrow\GL_4(\overline{\Q}_p),
 \qquad
 \rho_\pi^c\simeq\rho_\pi^\vee\eps^{-3}.
\]
It is unramified at every finite place nonsplit over $F$.
At every split finite place $v=ww^c$,
\begin{equation}
 \iota\WD\bigl(\rho_\pi|_{G_{L_w}}\bigr)^{\mathrm{F\text{-}ss}}
 \preceq
 \operatorname{rec}_w
        \bigl(\pi_v\otimes|\det|_w^{-3/2}\bigr).
 \label{eq:four-dimensional-compatibility}
\end{equation}
The representation is de Rham above $p$, and its increasing
weights at either distinguished $w\mid p$ are
\begin{equation}
 (h_1,h_2,h_3,h_4)
  =(1-\lambda'_1,\,2-\lambda'_2,\,
    3-\lambda'_3,\,-\lambda'_4).
 \label{eq:hodge-weights}
\end{equation}
The gaps $h_3-h_2$ and $h_4-h_3$ tend to infinity with the
choices in \cref{eq:shifted-weight}.
\end{proposition}

\begin{proof}
Apply \cref{thm:racsd-galois-input} to the representations
$\Pi_\beta=\tau_\beta\chi^{\delta_\beta}$ supplied by
\cref{lem:transfer-algebraicity}.
For a factor of rank $d$, length $a$, and $s\in S_a$, put
\[
 t_s=s+\frac{d-4}{2},\qquad
 \eta_s=\chi^{-\delta}|\cdot|_L^{t_s}.
\]
At a distinguished complex embedding its exponents are
$t_s-\delta/2$ and $t_s+\delta/2$.  Both are integers:
$2t_s\equiv a+d-5\equiv\delta\pmod2$.
Thus $\eta_s$ is an algebraic Hecke character even when its two
displayed factors are not separately algebraic.  Let $r_p(\eta_s)$
be its Galois character and define
\begin{equation}
 \rho_\pi
 =\bigoplus_\beta\ \bigoplus_{s\in S_{a_\beta}}
   r_p(\Pi_\beta)\otimes r_p(\eta_{\beta,s}).
 \label{eq:four-dimensional-construction}
\end{equation}
Its dimension is $\sum_\beta d_\beta a_\beta=4$.
On the automorphic side, the normalization in each summand is
\[
 (\tau\chi^\delta)|\det|^{(1-d)/2}\eta_s
      =\tau|\det|^{s-3/2}.
\]
Equations \eqref{eq:arthur-block-wd} and
\eqref{eq:racsd-compatibility} now give
\cref{eq:four-dimensional-compatibility}.
The same construction at the unramified standard parameters in
\cref{prop:simultaneous-transfer}\textup{(iii)} gives unramifiedness:
the bound forces the monodromy to vanish, and the semisimplified
Weil inertia is trivial.  Twisting and direct sums preserve
the de Rham property.

For completeness, the polarization follows directly with the stated
normalization.  Since $\chi^c=\chi^{-1}$,
\[
 \eta_s^c\eta_{-s}=|\cdot|_L^{d-4}.
\]
Geometric reciprocity sends $|\cdot|_L$ to $\eps$.  Consequently
the summand $R_s=r_p(\Pi)r_p(\eta_s)$ satisfies
\[
 R_s^c
   \simeq R_{-s}^{\vee}\eps^{1-d+d-4}
   =R_{-s}^{\vee}\eps^{-3}.
\]
The set $S_a$ is invariant under $s\mapsto-s$, proving the
polarization of \cref{eq:four-dimensional-construction}.

If $u$ is a first infinitesimal coordinate of $\tau$ and
$k=u+s$ is the corresponding coordinate of the full parameter,
the associated summand has de Rham weight
\[
 \frac{d-1}{2}-u-\frac\delta2
      -\left(t_s-\frac\delta2\right)
     =\frac32-k.
\]
Substitution from \cref{eq:transfer-infinitesimal} proves
\cref{eq:hodge-weights}.  In particular, for
$\lambda'=(m+H,m+H,1+H;-1-H)$ these weights are
\[
 1-m-H,\quad 2-m-H,\quad 2-H,\quad 1+H,
\]
with the asserted gaps $m$ and $2H-1$.
\end{proof}

\subsection{Local flags from the two unit operators}
\label{subsec:local-flags}

All the operator rings in this section are the residual local summands
specified in \cref{prop:cusp-congruence}.  In particular the two
normalized operators $\widetilde U_2,\widetilde U_3$ are already units.
The normalization is the one on the integral ordinary lattice, namely
$N_2=2$ and $N_3=2-H$ in the shifted weight.  It is not a condition on
the reduction of $A_0$.

\begin{lemma}[The two Hodge--Newton contacts]
\label{lem:two-contacts}
At each distinguished $w\mid p$, every $\rho_\pi$ under consideration
has a $G_{L_w}$-stable flag of dimensions $2,3,4$.  Its last quotient is
\begin{equation}
 X_w=\eps^{-h_4}\operatorname{unr}(\widetilde U_3^{-1}),
 \label{eq:last-character}
\end{equation}
where $\operatorname{unr}(a)$ takes geometric Frobenius to $a$.
At the conjugate place it has a first subline of character
\begin{equation}
 Y_{\bar w}=X_w^{-c}\eps^{-3}.
 \label{eq:first-character}
\end{equation}
Under \cref{eq:hecke-congruence} their values tend, uniformly on local
Galois elements and coefficientwise in $t$, to
$x=\eps^{-1}$ and $y=\eps^{-2}$, respectively.
\end{lemma}

\begin{proof}
Positive finite-slope Jacquet restriction at $w$
\citep[Section~4.1]{Casselman1995} allocates two
unramified characters to the final two slots of the Levi
$\GL_2\times\GL_1\times\GL_1$.  Write $\alpha_3,\alpha_4$ for their
geometric Frobenius eigenvalues with the cohomological norm shift of
\cref{prop:galois-package}.  The half-modulus in normalized Jacquet
restriction and the integral normalization of the two operators give
\begin{equation}
 \widetilde U_2=\frac{p^{h_3+h_4}}{\alpha_3\alpha_4},
 \qquad
 \widetilde U_3=\frac{p^{h_4}}{\alpha_4}.
 \label{eq:unit-root-formulas}
\end{equation}
Thus $v_p(\alpha_3)=h_3$ and $v_p(\alpha_4)=h_4$.  The sum of the
other two slopes is $h_1+h_2$, since the whole potentially semistable
module is weakly admissible.  Its Newton polygon lies above its Hodge
polygon.  The sum of the lowest two slopes is at most the sum of the
other two slopes and at least $h_1+h_2$; hence equality holds.
The strict Hodge gap at position two forces a strict Newton gap there.
Similarly the lowest three slopes sum to $h_1+h_2+h_3$, with a strict
gap at position three.  In particular the two final roots are
uniquely selected by their individual slopes.

After a finite semistable extension, the low-slope spaces are stable
under inertia and monodromy.  On either such space its induced Hodge
number is at least the sum of the corresponding lowest Hodge
weights, whereas weak admissibility bounds it above by its Newton
number.  Both inequalities are equalities.  The subspaces are
therefore weakly admissible.  The Colmez--Fontaine equivalence and
its subobject correspondence
\citep[Theorem~A, Proposition~4.2 and the corollary following
Theorem~4.3]{ColmezFontaine2000}
give semistable subrepresentations over this finite extension.
The argument
also applies to the descended module: the slope spaces are canonical
and hence stable under descent.  All embeddings above the fixed
embedding have the same Hodge multiset, so the same inequalities hold
after this extension.

The last quotient is unramified in its Weil part, has de Rham weight
$h_4$, and its Frobenius value is $\alpha_4$.  This proves
\cref{eq:last-character}; the polarization proves
\cref{eq:first-character}.  Their values belong to the integral
operator ring.  To check uniform convergence on local elements,
recall that $h_4=1+H_i$ and that $H_i$ is divisible by
$(p-1)p^{M_i+c}$.  Hence
\[
 u^{H_i}\equiv1\pmod{p^{M_i}}\qquad(u\in\Z_p^\times)
\]
uniformly in $u$, after the fixed choice of $c$.  This controls the
cyclotomic factor on the whole local Galois group.  For the unramified
factor, continuous powers of the unit $\widetilde U_3^{-1}$ define
the character on the procyclic unramified quotient.  Since
$\lambda_i^0(\widetilde U_3)=1$ in the full truncated ring, every
such power also has image $1$ there.  Formula~\eqref{eq:last-character}
therefore gives $X_w\equiv\eps^{-1}$ at the working precision,
uniformly on local elements and coefficientwise in $t$; the
polarization gives $Y_{\bar w}\equiv\eps^{-2}$.  These congruences
hold on every sequence of local elements.
\end{proof}

\subsection{Binary projections and small monodromy}
\label{subsec:binary-projections}

Away from $p$, the required local conditions come from inertia on the
final binary allocation.  We first separate its Frobenius roots from
the head, and then use the construction of $\rho_\pi$ to control
actual monodromy on that allocation.

At each split $(2,2)$ control place $v$ the positive central expansion
is a unit.  Jacquet lifting therefore supplies an unramified
spherical final binary block.  Let $A_v(Z)$ be its monic Frobenius
polynomial, obtained from the determinant and degree-one binary
Hecke operators.  The full degree-four polynomial $P_v(Z)$ also has
coefficients in the integral operator ring generated by the split
good spherical Hecke operators.  Indeed, for the finitely many cusp
systems at a fixed stage, choose good split Frobenius elements tending
simultaneously to the chosen local Frobenius.  Their characteristic
coefficients converge in the finite integral matrix lattice.  The
good operator subalgebra is already semisimple, finite over $\cO$,
and closed in that lattice; the limits consequently belong to that
subalgebra itself.  No normalization of a Hecke order is used here.
The same approximation, applied to \cref{eq:hecke-congruence}, gives
the comparison polynomial at every local element.

Perform monic division of $P_v$ by $A_v$ in the operator ring and
write $C_v$ for the monic quadratic quotient.  For forming an
operator, the possible remainder can be ignored.  On every
characteristic-zero system $A_vC_v=P_v$.  Set
\begin{equation}
 \Delta_v=\prod_{d=-4}^{4}
  \operatorname{Res}_Z\bigl(A_v(Z),Q_v^{2d}C_v(Q_v^{-d}Z)\bigr),
 \qquad \Delta=\prod_v\Delta_v,
 \label{eq:separation-resultant}
\end{equation}
where the product is over the bounded number of binary control
places, including the moving place.  All powers of $Q_v$ are
$p$-adic units.

\begin{lemma}[Separation and inertia]
\label{lem:binary-inertia}
The augmentation values of $\Delta$ are nonzero with uniformly
bounded valuation.  On any characteristic-zero system where
$\Delta\ne0$, the Frobenius primary projection $P_v^{\mathrm{bin}}$
onto the roots of $A_v$ is defined and satisfies
\begin{equation}
  (\rho_\pi(\iota)-1)P_v^{\mathrm{bin}}=0
  \qquad(\iota\in I_{L_v}).
 \label{eq:binary-inertia-identity}
\end{equation}
The assertion concerns actual inertia, without a semisimplicity
assumption on Frobenius.
\end{lemma}

\begin{proof}
At the moving place $Q_v=q_i$ tends to $1$ and the head Frobenius
matrix tends to $V(\gamma)$, while the binary roots tend to $1$.
The invertibility of $V(\gamma)-1$ in \cref{prop:tame-primes} gives
the required bound on its nonzero valuation.  At a fixed potentially good
place the head roots have the elliptic purity size, namely
$Q_v^{3/2}$ after the cohomological twist, whereas the final roots
are $Q_v,Q_v^2$ up to the specified fixed unitary character.
The half-integral difference in weights prevents equality after an
integral power of $Q_v$.  At a ramified potentially multiplicative
place choose the Frobenius lift whose quadratic sign is negative;
its head roots cannot equal the positive final roots after such a
power.  These are finitely many fixed nonzero comparisons.  The
weight and tame specializations converge to them, so their valuations
are uniformly bounded.

The nonzero resultants exclude both coincidences and norm linkages
between the two allocations.  The segment classification of
\citet{BernsteinZelevinsky1977} and
\citet[Theorem~6.1 and Proposition~8.5]{Zelevinsky1980}
then factors the local
representation into the unlinked head and binary representations.
A segment crossing the two allocations would yield a forbidden ratio
$Q_v^d$ with $|d|\leq4$.  The binary Jacquet allocation is therefore
the whole binary representation; its prescribed hyperspecial
invariants make it spherical, with zero predicted monodromy.

Separation from the head still allows norm linkages within the binary
pair.  Since \cref{eq:four-dimensional-compatibility} gives only
monodromy domination, we use the direct-sum construction of
$\rho_\pi$ to exclude those links as well.  In each individual shifted
copy of a cuspidal factor $\tau$ in
\cref{eq:four-dimensional-construction}, a selected root has a
length-one local segment.  A longer segment either crosses the
forbidden allocation boundary or lies wholly in the binary part,
contradicting sphericity.  Two length-one unramified data inside one
unitary generic $\tau_v$ cannot differ by $Q_v$: their real exponents
lie in an interval of width strictly less than one
\citep[Theorem~7.5]{Tadic1986} and
\citep[Theorem~9.7]{Zelevinsky1980}.  Thus the actual monodromy cannot
link two selected roots within that copy.  It cannot link different shifted copies,
since the Galois construction is their direct sum; it cannot link to
the complement by the resultants.  This proves that actual monodromy
vanishes on the selected part even though
\cref{eq:four-dimensional-compatibility} only gives an upper bound.

The semisimple Weil inertia is finite.  Its fixed subspace is stable
under Frobenius, and its selected generalized Frobenius primary
spaces have inertia one and monodromy zero.  Monic coprime
factorization gives the projection onto their sum, including any
Frobenius Jordan blocks.  This proves
\cref{eq:binary-inertia-identity}.
\end{proof}

\begin{lemma}[The Steinberg control]
\label{lem:steinberg-control}
At a split place with the specified $(3,1)$ parahoric level,
$\rho_\pi|_I$ is unipotent and
\[
 \rank(\rho_\pi(\iota)-1)\leq1\qquad(\iota\in I).
\]
\end{lemma}

\begin{proof}
Iwahori invariants make semisimple Weil inertia trivial.  Write the
local Langlands parameter as unramified Steinberg blocks of lengths
$s_j$.  The standard induction of these blocks surjects onto the
local representation.  Compact invariants are exact in characteristic
zero, so the standard induction must have $(3,1)$ parahoric
invariants.  Restriction to the maximal compact and its double-coset
decomposition reduces this to invariants of each Steinberg factor
under the induced partial flag.  The finite Weyl sign representation
of a Steinberg factor has such invariants only when all its positions
are separated.  A single-step flag can therefore accommodate at most
one block of length two, all others having length one.  The predicted
monodromy has rank at most one and square zero.  The monodromy bound
in \cref{prop:galois-package} gives the same for the actual
representation; its inertia matrices are $\exp(t_p(\iota)N)$.
\end{proof}

\subsection{Removing the operator nilpotents}

The cusp systems now have the required local subspaces and inertia
identities.  The same separated Jacquet allocations also control the
nilpotents in the operator order.  We use this fact to descend the
congruence of \cref{prop:cusp-congruence} to the algebra detected by
characteristic-zero eigensystems, retaining every coefficient of the
truncated tame variable.

\begin{proposition}[Uniform nilpotent control]
\label{prop:nilpotent-control}
Let $T_i$ be the image of $T_i^0$ in its characteristic-zero cusp
eigensystems.  There is an integer $C$, independent of $i$, the
weights and $q_i$, such that
\begin{equation}
 \Delta^C\ker(T_i^0\longrightarrow T_i)=0.
 \label{eq:nilpotent-annihilation}
\end{equation}
After decreasing the precisions by a fixed constant,
\cref{eq:hecke-congruence} descends to
\begin{equation}
 \lambda_i:T_i\longrightarrow
       \cO/(p^{M_i})[t]/(t^{b_*}),\qquad M_i\longrightarrow\infty.
 \label{eq:reduced-hecke-congruence}
\end{equation}
\end{proposition}

\begin{proof}
Work first on one irreducible local representation over an
algebraically closed field of characteristic zero.  Realize it as a
subquotient of induction from its full supercuspidal support.  The
geometric lemma filters the relevant block Jacquet module by at most
$4!=24$ allocation terms
\citep[Lemma 2.12]{BernsteinZelevinsky1977}.
Taking the required compact Levi invariants is exact.  On the
resulting finite-dimensional space, taking a joint generalized
eigenspace for the commuting controlled operators is also exact.
These operations therefore select a filtration of the actual
operator module.

At $p$, the unit summand was chosen before taking any reduced
quotient.  On every eigensystem its final slopes are $h_3,h_4$ by
\cref{lem:two-contacts}, and the other slopes are at most $h_2$.
The two gaps and \cref{eq:unit-root-formulas} select a unique
allocation to the $(2,1,1)$ Levi.  On its surviving grade the two
final $\GL_1$ center operators are scalars.  At a binary control
place where $\Delta_v\ne0$, the determinant and the degree-one
spherical operator determine the full quadratic $A_v$.  Separation
from the head similarly selects one allocation.  The hyperspecial
invariants of the full unramified binary induction are
one-dimensional, so all the binary operators are scalars on that
grade, including when the two binary inducing characters coincide.

Repeated factors within the head create no second allocation:
permutations inside that block already belong to the Levi Weyl group
in the geometric-lemma double-coset indexing.  No supercuspidal
factor splits internally across an allocation.  Thus there is
exactly one nonzero grade in the selected generalized eigenspace.
It equals that grade, and the controlled operators act by actual
scalars there.  In particular there is no residual Jordan extension
on the separated locus.  The good spherical operators are already
scalar on each automorphic representation.

For comparison, on a joint generalized component where $\Delta$
has eigenvalue zero, it acts by zero on every associated allocation
grade and therefore lowers the filtration.  A fixed power kills
that component.  Tensoring the filtrations at the bounded number
$r$ of control places gives, for example, the coarse bound $24^r$.
On the complementary components the kernel of eigensystem detection
is zero by the preceding scalar-action argument.  This proves
\cref{eq:nilpotent-annihilation} as a characteristic-zero operator
identity.  The integral order $T_i^0$ is torsion-free, so the same
identity holds integrally.

Finally the constant coefficient of $\lambda_i^0(\Delta)$ has
uniformly bounded nonzero valuation by \cref{lem:binary-inertia},
and all of its other coefficients are integral.  Successive
coefficient comparison in the truncated polynomial ring shows that
multiplication by its fixed power can lose only a bounded amount of
$p$-adic precision.  More explicitly, inversion over $k[t]/(t^{b_*})$
uses at most $b_*$ powers of the inverse constant coefficient; the
resulting denominator bound depends only on $C,b_*$ and that
valuation bound.  Applying this to
\cref{eq:nilpotent-annihilation} kills the kernel modulo $p^{M_i-O(1)}$.
Renaming these still increasing precisions proves
\cref{eq:reduced-hecke-congruence}.
\end{proof}
